
\section{Related Work}
\label{Related Work}
\paragraph{LLM-based Human Simulation}
LLM agents have been increasingly adopted as effective proxies for humans in research fields such as sociology and economics~\citep{xu2024ai,NBERw31122,gao2023large}.
In general, the usage of LLM agents can be categorized into \textbf{\textit{individual-level}} and \textbf{\textit{society-level}} simulation.
For the \textit{individual-level}, LLM agents have been leveraged to simulate individual activities or interactions, such as  human participants in surveys~\citep{argyle2023out}, humans' responses in HCI~\citep{hamalainen2023evaluating} or psychological studies~\citep{dillion2023can}, human feedback to social engineering attacks~\citep{asfour2023harnessing}, real-world conflicts~\citep{shaikh2023rehearsal}, users in recommendation systems~\citep{wang2023recagent,zhang2023generative}. 
For the \textit{society-level}, recent works have utilized LLM agents to model  social institutions or societal phenomenon, including a small town environment~\citep{park2023generative}, elections~\citep{zhang2024electionsim}, social networks~\citep{gao2023s3}, 
social media~\citep{tornberg2023simulating,rossetti2024social}, large-scale social movement~\citep{mou2024unveiling}, societal-scale manipulation~\citep{touzel2024simulation}, misinformation evolution~\citep{liu2024tiny}, peer review systems~\citep{jin2024agentreview},
macroeconomic activities~\citep{li2023large}, 
and world wars~\citep{hua2023war}. However, the majority of prior studies rely on an assumption without sufficient validation that \textit{LLM agents  behave like humans}. 
In this work, we propose a new concept, \textit{behavioral alignment}, to characterize the capacity of LLMs to simulate human behavior and discover that LLMs, particularly GPT-4, can largely simulate human trust behavior.



\paragraph{LLMs Meet Game Theory} 
The intersection of LLMs and Game Theory has attracted growing attention. The motivation is generally two-fold. One line of work aims to \textbf{\textit{leverage Game Theory to better understand LLMs' strategic capabilities and social behaviors}}. For example, \citet{akata2023playing,fan2023can,brookins2023playing} studied LLMs' interactive behaviors in classical games such as the Iterated Prisoner’s Dilemma. \citet{wang2023avalon,lan2023llm,light2023text,shi2023cooperation} explored LLMs' deception-handling and team collaboration capabilities in the Avalon Game.~\citet{xu2023exploring} discovered the emergent behaviors of LLMs such as camouflage and confrontation in a communication game Werewolf.
\citet{guo2024economics} discovered that most LLMs can show certain level of rationality in Beauty Contest Games and Second Price Auctions. \citet{mukobi2023welfare} measured the cooperative capabilities of LLMs in a general-sum variant of Diplomacy.
\citet{guo2023suspicion} proposed to elicit the theory of
mind (ToM) ability of GPT-4 to play various imperfect information games.
The other line of works aims to \textbf{\textit{study whether or not LLM agents can replicate existing human studies in Game Theory}}. This direction is still in the initial stage and needs more efforts. One typical example is \citep{aher2023using}, which attempted to replicate existing findings in studies such as the Ultimatum Game. Another recent work explored the similarities and differences between humans and LLM agents regarding emotion and belief in ethical dilemmas~\citep{lei2024fairmindsim}.
Different from previous works, we focus on a critical but under-explored behavior, \textit{trust}, in this paper and reveal it on LLM agents. 
We also discover the \textit{behavioral alignment} between agent trust and human trust with evidence in both \textit{actions} and \textit{underlying reasoning processes}, which is particularly high for GPT-4, implying that LLM agents can not only replicate human studies but also align with humans' underlying reasoning paradigm. Our discoveries illustrate the great potential to simulate human trust  behavior with LLM agents.























